\section{Discharging $\AND$: Hadamard as Ideal Membership}
\label{sec:hadamard}

This section proves the nonlinear part of the UAIR: the column Hadamard
products that encode $\AND$ (\Cref{sec:and}) and the carry Hadamards of the
modular additions (\Cref{sec:add}). The idea is to make the nonlinearity look
like everything else. Reinterpreting each cell's bits as the coefficients of a
univariate over a small subspace turns a bitwise $\AND$ into a polynomial
\emph{ideal-membership} claim, of exactly the kind the linear part already
discharges (\Cref{sec:piop}). The Hadamard relations then ride the same
machinery --- one ideal check, the evaluation projection $\psia$, and the
Step-4 sumcheck --- at the \emph{same} challenge point, settled by the
\emph{same} commitment opening. The only nonlinearity in SHA-256 thus costs the
prover one extra polynomial message and one extra group in a sumcheck it is
already running, and costs the commitment nothing beyond reading a second
functional off the one bound object.

\subsection{The obligation}
\label{sec:had-oblig}

Collect the $K_H$ Hadamard relations of the trace --- the two $\AND$
products per round ($\Ch$ and $\Maj$) and the carry Hadamard of each binary
addition step, $K_H = 14$ in the deployed layout --- as triples of operand
columns
\[
  \bp{u}^{(k)}\had\bp{v}^{(k)} \;=\; \bp{w}^{(k)},
  \qquad k = 1,\dots,K_H,
\]
where each operand is a fixed $\Ftwo$-linear combination of committed
columns and their shifts: for the $\AND$ relations, committed columns and
their \emph{row}-shifts (e.g.\ $\col{E}$ and $\col{E}^{\downarrow1}$); for
the adder relations, the operands $\bp{a}+X\bp{c}$, $\bp{b}+X\bp{c}$,
$\bp{c}+X\bp{c}$ of \Cref{lem:binius}, in which $X\bp{c}=\SHL^{1}(\bp{c})$ is a
\emph{within-cell bit-shift} of the carry, not a row-shift. The relation \eqref{eq:had-ideal} below is a
coefficient-wise identity among these operands as bit-vectors, so the
discharge is indifferent to the distinction; it resurfaces only at binding
(\Cref{sec:had-binding},~\Cref{rem:adder-trusted}). Writing
$u^{(k)}_b(x)\in\{0,1\}$ for bit $b$ of the operand at row
$x\in\{0,1\}^{\nu}$, the obligation is the system of scalar equations
\begin{equation}
  \label{eq:had-system}
  u^{(k)}_b(x)\,v^{(k)}_b(x) \;=\; w^{(k)}_b(x)
  \qquad \text{for all } k,\; b\in\{0,\dots,\wbit-1\},\; x\in\{0,1\}^{\nu}.
\end{equation}

\subsection{Bits as Lagrange coefficients: $\AND$ as ideal membership}
\label{sec:psi-z}

Identify the $\wbit$ bit positions with an $\Ftwo$-affine subspace
$H\subset\K$ of dimension $\mu=\log_2\wbit$, and let $\{L_b\}$ be the Lagrange
basis of $H$. Read a cell's bit-vector $\bp{w}$ not as monomial coefficients
(the reading $\psia$ uses) but as the \emph{Lagrange} coefficients of the
degree-$<\wbit$ univariate
\[
  L(\bp{w})(X) \;=\; \sum_b w_b\,L_b(X),
  \qquad L(\bp{w})(h_b) = w_b \ \text{ for the base points } h_b\in H .
\]
The point of this reading is a clean algebraic restatement of $\AND$. Since
$L(\bp{u}),L(\bp{v}),L(\bp{w})$ take the witness bits on the base domain $H$,
the per-bit products $u_bv_b=w_b$ of \eqref{eq:had-system} hold for all $b$
exactly when the polynomial $L(\bp{u})\cdot L(\bp{v}) - L(\bp{w})$ vanishes on
all of $H$ --- that is, exactly when it lies in the ideal generated by
$Z_H(X)=\prod_{h\in H}(X-h)$:
\begin{equation}
  \label{eq:had-ideal}
  \bp{u}^{(k)}\had\bp{v}^{(k)} = \bp{w}^{(k)}
  \quad\Longleftrightarrow\quad
  L(\bp{u}^{(k)})\,L(\bp{v}^{(k)}) - L(\bp{w}^{(k)}) \;\in\; \langle Z_H\rangle .
\end{equation}
This is the same shape as the linear part's obligations (\Cref{sec:piop}),
which are membership in $\langle 0\rangle$ or $\langle X\rangle$; the
Hadamard part simply adds the generator $Z_H$. Batching the $K_H$ relations by a transcript challenge
$\gamma_{\mathrm{rel}}$ and the rows by the usual $\eq(\cdot;r)$, the whole
nonlinear obligation is the single ideal-membership claim
\begin{equation}
  \label{eq:had-zerocheck}
  \sum_{x\in\{0,1\}^{\nu}} \eq(x;r)
  \sum_{k} \gamma_{\mathrm{rel}}^{\,k}
  \Bigl(L(\bp{u}^{(k)})\,L(\bp{v}^{(k)}) - L(\bp{w}^{(k)})\Bigr)(X,x)
  \;\in\; \langle Z_H\rangle ,
\end{equation}
where $L(\bp{u}^{(k)})(X,x)=\sum_b u^{(k)}_b(x)L_b(X)$ is the operand's
Lagrange univariate at row $x$.

\subsection{Discharging it on the main pipeline}
\label{sec:had-pipeline}

Membership \eqref{eq:had-zerocheck} is checked by the same three steps as the
linear ideal check, and at the same points.

\paragraph{The ideal-check message.} The prover sends the witness for
membership in $\langle Z_H\rangle$: the off-subspace values of
\[
  R_0(X) \;=\; \sum_{x}\eq(x;r)\sum_k\gamma_{\mathrm{rel}}^{\,k}
  \bigl(L(\bp{u}^{(k)})\,L(\bp{v}^{(k)}) - L(\bp{w}^{(k)})\bigr)(X,x),
  \qquad \deg R_0 \le 2(\wbit-1).
\]
Two economies make this cheap. On the \emph{base} domain $X\in H$ the operands
take the $\{0,1\}$ witness bits and $R_0$ \emph{vanishes for honest data} by
\eqref{eq:had-ideal} --- so the prover transmits only the $\wbit-1$
\emph{off-subspace} values, the genuine content of the $\langle Z_H\rangle$
witness. And the dominant accumulation over the $2^{\nu}$ rows, restricted to
the base domain, is $\Ftwo$ (bit) arithmetic; only the off-subspace lift
touches $\K$. The message is absorbed into the transcript before $\alpha$, as
the linear ideal-check messages are.

\paragraph{The projection at $\alpha$.} The same evaluation point $\alpha$
that projects the linear constraints settles membership: a degree-$<\wbit$
polynomial lies in $\langle Z_H\rangle$ iff it vanishes on $H$, and the
verifier tests this at $X=\alpha$ via the shipped off-subspace values. Reading
the operands at $\alpha$ is the second instance of the cell projection of
\Cref{def:projection},
\[
  \psiz(\bp{w}) \;=\; L(\bp{w})(\alpha) \;=\; \sum_b L_b(\alpha)\,w_b \;\in\;\K ,
\]
the subspace-Lagrange weighting $\xi_b=L_b(\alpha)$ --- the exact analogue of
$\psia$'s monomial weighting $\xi_b=\alpha^b$, read at the \emph{same} point
$\alpha$. The product term then joins the Step-4 sumcheck
(\Cref{sec:sumcheck}) as one additional degree-$3$ group,
\[
  \sum_{x}\eq(x;r)\sum_k\gamma_{\mathrm{rel}}^{\,k}
  \bigl(\psiz(\bp{u}^{(k)})(x)\,\psiz(\bp{v}^{(k)})(x) - \psiz(\bp{w}^{(k)})(x)\bigr),
\]
whose claimed sum is $R_0(\alpha)$, reconstructed by the verifier from the
off-subspace message. The bit dimension never becomes a sumcheck variable: the
$\mu$ bit positions are folded once, at $\alpha$, by the Lagrange reading.

\paragraph{The shared sumcheck.} Because this group rides the Step-4 sumcheck
that already reduces the linear constraints, it is reduced over the same $\nu$
row variables to the same point $r^*$, yielding the closing operand
evaluations
\begin{equation}
  \label{eq:operand-evals}
  \psiz\bigl(\bp{u}^{(k)}\bigr)(r^*),\quad
  \psiz\bigl(\bp{v}^{(k)}\bigr)(r^*),\quad
  \psiz\bigl(\bp{w}^{(k)}\bigr)(r^*),
\end{equation}
whose combination $\sum_k\eq(k;\gamma_{\mathrm{rel}})\,(a^{(k)}b^{(k)}-c^{(k)})$
the verifier checks against the group's reduced claim. There is no second
sumcheck and no second challenge point: the nonlinear part shares the linear
part's $r^*$.

\begin{lemma}[The Hadamard reduction is sound]
\label{lem:psi-z-sound}
If \eqref{eq:had-system} holds then \eqref{eq:had-zerocheck} holds and $R_0$
vanishes on $H$, so the check passes for every $\gamma_{\mathrm{rel}},r$ and
every $\alpha$. Conversely, if some equation of \eqref{eq:had-system} fails,
then over uniform $\gamma_{\mathrm{rel}}, r, \alpha$ the check passes with
probability at most
$\frac{2(\wbit-1)}{|\K|}+\frac{K_H-1}{|\K|}+\frac{\nu}{|\K|}$.
\end{lemma}

\begin{proof}[Idea]
A failed equation makes the bit-residual
$\bp{u}^{(k)}\had\bp{v}^{(k)}-\bp{w}^{(k)}$ nonzero at some base point, so
$L(\bp{u}^{(k)})L(\bp{v}^{(k)})-L(\bp{w}^{(k)})$ does \emph{not} vanish on $H$
and $R_0$ is a nonzero univariate of degree $\le 2(\wbit-1)$; it is therefore
nonzero at the random $\alpha$ except with probability
$\frac{2(\wbit-1)}{|\K|}$ (Schwartz--Zippel). The relation batch by
$\gamma_{\mathrm{rel}}$ and the $\eq(\cdot;r)$ row check add the remaining two
terms by the standard arguments. The dominant term is the analogue of the
linear ideal check's $\frac{\wbit-1}{|\K|}$; $\psia$ and $\psiz$ are the same
evaluation at $\alpha$, read in the monomial resp.\ Lagrange basis.
\end{proof}

The reason the \emph{Lagrange} basis is needed at all is
\Cref{rem:and-not-survive}: it keeps the bit-products
$u_bv_b$ \emph{separated by bit position}, so the verifier's final check is a
genuine product of per-position bits, whereas $\psia$ of $\bp{u}\had\bp{v}$
would scramble them into cross-terms. Both readings, however, live at the one
point $\alpha$ and are read off one committed object, as the binding now shows.

\begin{remark}[The $\mathrm{GF}(2^8)$ acceleration]
\label{rem:gf8}
The base-domain accumulation can be pushed below even $\Ftwo$-vs-$\K$ by
running the additive NTT that evaluates $R_0$ over an \emph{embedded}
$\mathrm{GF}(2^8)$ subspace with a byte-lookup table, rather than over the
monomial subspace of $\K$ directly. This is a pure performance variant: the
embedding $\mathrm{GF}(2^8)\hookrightarrow\K$ is a field homomorphism, so the
message computed in $\mathrm{GF}(2^8)$ equals its image in $\K$, and the
soundness of \Cref{lem:psi-z-sound} is unchanged. In the deployed prover it is
several-fold faster than the monomial path and carries the bulk of the
Hadamard-discharge speedup.
\end{remark}

\subsection{Binding the discharge to the commitment}
\label{sec:had-binding}

The reduction has left the operand evaluations \eqref{eq:operand-evals} at
$r^*$. Two facts let these ride the \emph{same} opening that settles the linear
part, with no second commitment and no second point.

\emph{$\AND$ operands are $\Ftwo$-linear in column row-shifts.} Each $\AND$
operand is an $\XOR$ of committed columns and their row-shifts, and $\psiz$
is $\Ftwo$-linear (\Cref{def:projection}), so $\psiz(\bp{u}^{(k)})(r^*)$ is
the $\Ftwo$-sum of the per-column \emph{pair-evals}
$\psiz(\col{c}^{\downarrow\Delta})(r^*)$ of its constituent columns ---
exactly as the linear path needs the $\psia$ pair-evals, and through the same
row-shift machinery. (The adder operands also carry a within-cell bit-shift
$X\bp{c}=\SHL^{1}(\bp{c})$, bound by shifted Lagrange weights rather than by a
row-shift; see \Cref{rem:adder-trusted}.)

\emph{Pair-evals fold into the multipoint eval.} The
$\psiz(\col{c}^{\downarrow\Delta})(r^*)$ are appended to the main multipoint
evaluation (\Cref{sec:multipoint}) as pointed-shift claims and reduced ---
alongside the $\psia$ claims, to which they are wholly analogous --- to the
single point $r_0$. There the commitment opening binds, for each column, the
entire bit-slice fold $a'_g$ (a degree-$<\wbit$ polynomial over $\K$ whose
$b$-th coefficient is the $b$-th bit slice folded to $r_0$; see
\Cref{sec:commitment}). Because one bound object carries all bit slices, and
because both readings live at the same point $\alpha$, the verifier reads
\emph{both} weightings off it:
\[
  \psia(a'_g) = \mle{\psia(w_g)}(r_0)\ \ (\text{linear claim}),
  \qquad
  \psiz(a'_g) = \mle{\psiz(w_g)}(r_0)\ \ (\text{discharge claim}).
\]
Concretely, the open exposes its column-batching challenge $\gamma$, and the
single binding check
\begin{equation}
  \label{eq:psi-z-binding}
  \psiz(a') \;\stackrel{?}{=}\; \sum_{g} \gamma_g\cdot\mle{\psiz(w_g)}(r_0),
  \qquad a' = \sum_g\gamma_g a'_g,
\end{equation}
ties the discharge's evaluations to the committed bit slices. A final
\emph{tie} recombines the now-bound pair-evals into the operand evaluations
\eqref{eq:operand-evals} (by the $\Ftwo$-linearity above) and checks them
against the group's closing claim. This is the entire cost of binding: one
extra weight-vector projection \eqref{eq:psi-z-binding} of the already-bound
$a'$, no new opening and no new point.

\begin{remark}[Adder operands are trusted]
\label{rem:adder-trusted}
The $\AND$ operands are $\XOR$s of committed columns and their row-shifts, so
they bind as above. The adder operands look more delicate only because of the
term $X\bp{c}$ --- but $X\bp{c}=\SHL^{1}(\bp{c})$ is just a \emph{within-cell
shift} of the carry (\Cref{lem:binius}), not an algebraic factor that must be
``moved across'' the $\langle Z_H\rangle$ discharge. Two of the three operands
are not even shifts: the linear identity \eqref{eq:add-linear} with its bit-$0$
check forces $X\bp{c}=\bp{D}$ (with $\bp{D}=\bp{s}+\bp{a}+\bp{b}$), so
\[
  \bp{a}+X\bp{c} = \bp{s}+\bp{b},
  \qquad
  \bp{b}+X\bp{c} = \bp{s}+\bp{a}
\]
are plain $\XOR$s of committed columns, binding exactly like an $\AND$ operand.
Only $\bp{c}+X\bp{c} = \SHR^{1}(\bp{D})+\bp{D}$, with the carry-out bit
$\beta$, keeps a within-cell shift --- and a within-cell shift binds the same
way a row-shift does, one dimension over: by reading the bound bit-slice fold
with \emph{shifted} Lagrange weights,
$\psiz(\SHL^{1}\bp{c}) = \sum_b L_{b+1}(\alpha)\,a'_b$, still a single
length-$\wbit$ read-off of $a'$. The deployed prover does not yet perform this
bit-shift read-off and does not commit $\beta$, so it ships the adder $\psiz$
parents \emph{trusted} (honest-prover-supplied) and folds them into the tie.
Adding the bit-shift binding and committing $\beta$ closes the gap; it is
localized to the adder operands and leaves the $\AND$ relations and the entire
linear part unconditional. A recorded follow-up.
\end{remark}

\subsection{Summary}
\label{sec:had-summary}

The Hadamard part of SHA-256 --- its only nonlinearity --- is an
ideal-membership claim \eqref{eq:had-ideal} in $\langle Z_H\rangle$, obtained
by reading cell bits as Lagrange coefficients. It is discharged by the linear
part's own machinery: one ideal-check message (the off-subspace values of
$R_0$), one extra degree-$3$ group in the Step-4 sumcheck at the shared point
$\alpha$, reduced to the shared point $r^*$, with soundness
$O\!\bigl((\wbit + K_H + \nu)/|\K|\bigr)$ (\Cref{lem:psi-z-sound}). It binds to
the commitment by the single projection check \eqref{eq:psi-z-binding} on the
shared bit-slice opening. The discharge thus costs, at the commitment layer,
nothing beyond reading the cell in a second (Lagrange) basis off the one object the opening
already binds --- the property the next section's commitment is built to
provide.
